\documentclass[10pt,a4paper]{amsart}
\usepackage[margin=19mm]{geometry}
\usepackage{amsmath,amssymb,mathtools}
\usepackage[scr=rsfs,cal=euler]{mathalfa}
\usepackage{xfrac}
\usepackage[unicode,hidelinks,bookmarks=false]{hyperref}
\newcommand{\ie}{i.e.\ }
\newcommand{\cf}{cf.\ }
\newcommand{\resp}{respectively\ }
\newcommand{\Subsection}{\subsection}
\providecommand{\nobreakdash}{\nobreak\hbox{-}}
\setlength{\emergencystretch}{2em}
\multlinegap0pt
\usepackage{amssymb}
\usepackage{mathtools}
\newtheorem{thm}{Theorem}
\title{Two $t$-analogues of the\\ tree inversion enumerator}
\author{Sam Hopkins}
\date{Source: arXiv:2510.22385v2 (CC BY 4.0)}
\begin{document}
\maketitle

\noindent\emph{Source and licence.} Two $t$-analogues of the\\ tree inversion enumerator. Version arXiv:2510.22385v2 (CC BY 4.0), distributed by arXiv under the Creative Commons Attribution 4.0 International licence, http://creativecommons.org/licenses/by/4.0/, as recorded in the arXiv licence metadata for this exact version and shown on the abstract page https://arxiv.org/abs/2510.22385v2. Full licence notice: \texttt{licenses/arxiv-2510.22385-CC-BY-4.0.txt}.\par\smallskip

\noindent\textit{Editorial scope.} Counted unit: the article's thm thm (source0.tex lines 60--65). The article's complete original proof of it is delivered with it (source0.tex lines 66--74, one \texttt{proof} environment of 9 lines). No mathematics is written, repaired or re-derived: the statement and the proof are the article's own lines, in the article's own notation, and no retained line is substituted or re-flowed.  No further source-local context is needed: every cross-reference of every retained line resolves inside the two delivered spans.  Nothing was omitted from any retained span, so the package carries no omission record.  No retained line contains a citation, so the wrapper carries no bibliography and no bibliography entry.  No retained line contains a figure environment, an image-inclusion command or drawing code, so no image is delivered and the package declares no asset dependency.  Editorial formatting only, in addition: an AMS \texttt{amsart} preamble that restates the article's own declarations and macros used by the retained lines, namely the environments \texttt{thm} on separate counters, together with the standard packages those lines need.  The retained environments print in this document as Theorem 1 in order of appearance, whereas the article prints the same items with the numbering of its own full document.  The complete native source of the article as distributed in the arXiv e-print for this exact version is bundled unchanged as \texttt{sources/187794/source0.tex}.
\begin{thm}
We have
\[I_n(q,0)=\widetilde{I}_n(q,0)= [n]_q!\]
and
\[I_n(0,t)=\widetilde{I}_n(0,t) = A_n(t).\]
\end{thm}

\begin{proof}
To show $I_n(q,0)=[n]_q!$: this is clear because trees with one leaf are permutations in a natural way, and the tree inversions are inversions of the permutation.

To show $\widetilde{I}_n(q,0) = [n]_q!$: parking functions $\pi\in \mathrm{PF}(n)$ with $\mathrm{exced}(\pi)=0$ are sequences of integers $(\pi_1,\pi_2,\ldots,\pi_n)$ with $1 \leq \pi_i \leq i$ for all $i$. (These are essentially Lehmer codes of permutations.) The cosum enumerator of these is evidently $[n]_q!$.

To show $I_n(0,t)=A_n(t)$: it is well-known that trees $T \in \mathrm{Tree}(n+1)$ with $\mathrm{inv}(T)=0$, i.e., \emph{increasing trees}, are also in bijection with permutations in $S_n$. Indeed, given such an increasing tree, perform a depth-first search on the tree, starting at the root $0$, and always preferring to visit the vertex with the biggest label when one has a choice. Record the order in which vertices are visited in this search to obtain a permutation. This is a bijective procedure. Moreover, the leaves of the tree correspond to the descents of the permutation, except for the ``last'' leaf.

To show $\widetilde{I}_n(0,t)=A_n(t)$: parking functions with cosum zero are permutations, and it is well-known that the generating function of permutations by excedances is the Eulerian polynomial.
\end{proof}

\end{document}
